\begin{definition}[Banach space over a complete valued field] \label{definition:banach_space_over_a_commutative_ring_with_an_absolute_value_valued_in_an_ordered_semiring}
    \TODO{AI generated, verify}
    Let $T$ be an \CrefAndHyperrefIfExist{definition:ordered_semiring}{ordered semiring}, let $R$ be a \CrefAndHyperrefIfExist{definition:commutative_ring}{commutative ring} equipped with an \CrefAndHyperrefIfExist{definition:absolute_value_on_a_ring_valued_in_an_ordered_semiring}{absolute value valued in $T$}:
    $$|\cdot| : R \to T,$$
    such that $R$ is \CrefAndHyperrefIfExist{definition:complete_extended_metric_space}{complete} with respect to the \CrefAndHyperrefIfExist{definition:metric_induced_by_an_absolute_value_on_a_commutative_ring}{metric induced by} $|\cdot|$.

    A \hldef{Banach space over $R$ valued in $T$} or simply a \hldef{complete extended normed module over $R$ valued in $T$} is a \CrefAndHyperrefIfExist{definition:module_over_a_ring}{module} $M$ over $R$ equipped with an \CrefAndHyperrefIfExist{definition:extended_norm_on_a_module_over_a_commutative_ring_with_absolute_value}{extended norm valued in $T^\infty$}:
    $$\|\cdot\| : M \to T,$$
    (taking values in $T$, not $\infty$) such that the metric $d(x,y) := \|x - y\|$ \CrefAndHyperrefIfExist{definition:extended_metric_induced_by_an_extended_norm_on_a_module_over_a_commutative_ring_with_absolute_value}{induced by the norm} makes $(M,d)$ into a \CrefAndHyperrefIfExist{definition:complete_extended_metric_space}{complete extended metric space}.

    Two conditions on $T$ ensure that completeness is non-vacuous and well-behaved. First, if $\{\varepsilon \in T : \varepsilon > 0\}$ has a minimum element $\varepsilon_0$, then every Cauchy sequence $(x_n)$ eventually satisfies $d(x_n,x_m) < \varepsilon_0$, which forces $d(x_n,x_m) = 0$ and the sequence becomes eventually constant; completeness holds trivially. More generally, if for some $\delta > 0$ there is no $\varepsilon \in T$ with $0 < \varepsilon < \delta$, then any Cauchy sequence eventually has pairwise distances strictly less than $\delta$, which again forces eventual constancy. Thus non-vacuous completeness requires $T$ to be \hl{dense at zero} (for every $\delta > 0$ there exists $0 < \varepsilon < \delta$). Second, if the positive cone is not downward-directed, then the Cauchy condition "for every $\varepsilon > 0$" can be satisfied for some positive elements while failing for others that share no common lower bound, so it does not correspond to any genuine notion of convergence. Total order implies both properties (minimum exists iff minimal exists, and downward-directedness holds by taking minima). The standard case $T = \bbR_{\geq 0}$ satisfies neither pathology: its positive cone is totally ordered, downward-directed, and dense at zero, ensuring meaningful completeness.

    In the literature, "Banach space" usually means the special case where $R = k$ is a \CrefAndHyperrefIfExist{definition:field}{field}, $M$ is a \CrefAndHyperrefIfExist{definition:vector_space_over_a_field}{vector space} over $k$, and $T = \bbR_{\geq 0}$. If the underlying valued field $(k, |\cdot|)$ is \CrefAndHyperrefIfExist{definition:archimedean_absolute_value_on_a_field}{non-archimedean}, then one often speaks of a "\hldef{non-archimedean Banach space}" and sometimes additionally assumes that the norm satisfies the ultrametric inequality $\|x + y\| \leq \max(\|x\|, \|y\|)$.
\end{definition}